\documentclass[11pt,a4paper]{article}
\usepackage[T1]{fontenc}
\usepackage[utf8]{inputenc}
\usepackage{lmodern}
\usepackage[margin=22mm]{geometry}
\usepackage{microtype}
\usepackage{array,booktabs,longtable}
\usepackage{enumitem}
\usepackage{xcolor}
\usepackage{amsmath}
\usepackage[obeyspaces]{url}
\usepackage{hyperref}
\usepackage{xurl}
\usepackage{needspace}
\hypersetup{colorlinks=true,linkcolor=blue!45!black,urlcolor=blue!45!black,
  pdftitle={ProcRosetta: A Simple Guide to the Streamlit Interface},
  pdfauthor={ProcRosetta interface documentation}}
\setlength{\parindent}{0pt}
\setlength{\parskip}{5pt}
\setlength{\emergencystretch}{2em}
\setlist{nosep,leftmargin=*}
\AddToHook{cmd/section/before}{\Needspace{12\baselineskip}}
\AddToHook{cmd/subsection/before}{\Needspace{10\baselineskip}}
\renewcommand{\arraystretch}{1.15}
\newcommand{\ui}[1]{\textbf{#1}}
\newcommand{\code}[1]{\nolinkurl{#1}}
\newenvironment{controls}{%
  \small\begin{longtable}{@{}>{\raggedright\arraybackslash}p{.34\linewidth}
    >{\raggedright\arraybackslash}p{.62\linewidth}@{}}
  \toprule\textbf{On the screen} & \textbf{What it means or does}\\\midrule\endhead
  \bottomrule\endfoot
}{\end{longtable}\normalsize}
\title{ProcRosetta\\A Simple Guide to the Streamlit Interface}
\author{Standalone user guide}
\date{Interface reviewed: 25 September 2026}

\begin{document}
\maketitle

This guide explains the pages, buttons, settings, tables, and diagrams in the
ProcRosetta Streamlit studio. You can use it without knowing how the neural
model is implemented. Names in \ui{bold} match labels in the interface.

The current application has a \ui{Studio overview} and four main views:
\ui{Artifact workspace}, \ui{Translation studio}, \ui{Latent explorer}, and
\ui{Checkpoint dashboard}. A separate \ui{Quality evaluation} page exists in
the repository but is not included in the application's navigation. It is
covered in Appendix~\ref{sec:evaluation} so that the guide also explains that
interface without suggesting that it is available in the normal menu.

This is an independent document. It does not need to be included in the main
paper. Its descriptions follow the source files listed in
Appendix~\ref{sec:sources}.

\tableofcontents
\newpage

\section{Start here}
\label{sec:start}

\subsection{What the studio does}
The studio takes a description of a process, turns it into a list of numbers,
and uses those numbers to compare processes or generate a process model.
There are three kinds of input:

\begin{controls}
Event log (\code{.xes}) & A record of what happened. One \emph{trace} is one
case, such as one order. Its events list the activities performed in order.\\
Process tree (\code{.ptml}) & A model that combines activities using sequence,
choice, parallel work, and repetition.\\
Petri net (\code{.pnml}) & A model with places, transitions, and connecting
arrows. Tokens in places describe the process state.\\
\end{controls}

The interface calls an imported file an \emph{artifact}, and its type a
\emph{modality}. An \emph{embedding} or \emph{latent vector} is the list of
numbers describing an artifact. \emph{Encoding} creates that list.
\emph{Decoding} turns a list back into a process tree. A \emph{checkpoint}
is a saved trained model used for both steps.

The usual path through the application is:
\begin{center}
\fbox{\begin{minipage}{.92\linewidth}\centering
Choose a checkpoint $\longrightarrow$ import a file $\longrightarrow$
inspect its input\\[3pt]
$\longrightarrow$ compare embeddings or decode a tree
$\longrightarrow$ inspect and download the result
\end{minipage}}
\end{center}

The tree is the model's direct output. The displayed output Petri net is
converted from that tree. A valid output can still describe different behavior
from the input, so check the source comparisons as well as the validation badges.

\subsection{A first walkthrough}
\begin{enumerate}
\item Start the application from the repository root with
  \code{streamlit run streamlit_app.py}, and open the browser address printed
  by Streamlit. The dependencies must already be installed.
\item In the sidebar, choose a \ui{Checkpoint}. A trusted \code{.pt} file
  must be installed in the server's checkpoint directory.
\item Open \ui{Artifact workspace}. In \ui{Import artifacts}, select
  \code{running-example.xes} under \ui{Bundled artifacts}, then press
  \ui{Import bundled selection}. Importing also encodes the artifact.
\item Select it under \ui{Inspect artifact}. Read \ui{Source preview},
  \ui{Model-facing input}, and \ui{Warnings \& state}.
\item Open \ui{Translation studio}, select the log under \ui{Latent sources},
  keep \ui{Deterministic mean}, and press \ui{Decode process tree}.
\item Inspect \ui{Validation gates}, \ui{Decoded process tree}, and
  \ui{Source comparison}. Use \ui{Download PTML} to save a generated tree
  when that button is available.
\item To compare different representations, also import
  \code{running-example.ptml} and \code{running-example.pnml}. Bundled files
  with the same base name receive the same process group. Open
  \ui{Latent explorer} to compare their embeddings.
\end{enumerate}

\section{Navigation, the sidebar, and shared behavior}
\label{sec:shared}

\subsection{Finding a page}
\begin{controls}
Studio overview & Starting page, workspace counts, and shortcuts to the four views.\\
Artifact workspace & Import, inspect, group, encode, remove, and export artifacts.\\
Translation studio & Decode embeddings into trees, inspect the generation steps,
compare with sources, and download outputs.\\
Latent explorer & Compare embeddings using plots, distances, neighbors, and
experimental interpolation.\\
Checkpoint dashboard & Inspect the saved model, training history, configuration,
and other trusted checkpoints.\\
\end{controls}

Use the page links in the sidebar to move between views. If the sidebar is
collapsed, expand it to see the navigation and model settings. The workspace is
shared between pages in the same Streamlit session. It is held in session
memory; save useful results before refreshing into a new session or closing
your work. There is no workspace-ZIP restore button.

Tabs switch between related panels. Expanders, such as \ui{Checkpoint facts},
open and close extra details. A checkbox enables an option; a radio control
selects one option; a multiselect allows several choices. Number fields and
sliders change a setting. Buttons carry out an action. Greyed-out buttons need
some prerequisite, such as an imported artifact or a selected source.

Tables may need horizontal scrolling to reveal all columns. Hover over a
chart to inspect its values where tooltips are available. Progress bars and
loading messages indicate work in progress. Green messages report success,
yellow messages explain limitations, red messages report errors, and information
messages often explain why a page is empty.

\subsection{Active checkpoint}
\begin{controls}
Checkpoint & Select a saved model by filename. Only trusted files already on
the server are listed; the browser has no checkpoint-upload control.\\
Inference device & Choose where model calculations run. CPU is always offered;
the detected default accelerator is also offered when available.\\
Epoch / latent / activities caption & The saved training epoch, the number of
coordinates in each embedding, and the model's activity-vocabulary limit.\\
Checkpoint facts & Opens structured text (JSON) containing the identifier,
checkpoint type, epoch, latent and hidden dimensions, maximum activities,
maximum tree arity, and best validation loss. The identifier distinguishes
models; arity means the number of children of a tree operator.\\
\end{controls}

The server reads checkpoints from \code{checkpoints/}, or from the directory
specified by \code{PROC_ROSETTA_CHECKPOINT_DIR}. A missing or unreadable
checkpoint produces an error. The overview remains visible, but the four
working views stop before their main controls; artifact import also needs a
loadable checkpoint in this interface.

Changing the checkpoint does not automatically re-encode existing artifacts.
Old entries can show \code{checkpoint mismatch}; Translation studio and Latent
explorer use only compatible embeddings. Import the files again with the new
checkpoint to encode them for that model. Existing register counts can still
include stored results from an earlier checkpoint. Use one checkpoint for all
items before making a comparison or preparing a workspace export.

\subsection{When settings take effect}
File selection starts import automatically. Process-group edits take effect
immediately. Plot controls update their displays. Decoding, interpolation,
gallery building, evaluation, and export preparation require their action
buttons. Changing one of their settings does not replace an already displayed
result: run the action again before interpreting or downloading it.

\section{Studio overview}
The heading is \ui{See every transformation.} The four counters describe the
current workspace:

\begin{controls}
Workspace artifacts & Number of imported artifacts.\\
Embeddings ready & Artifacts with a stored, nonempty embedding and no encoding
errors. This count does not filter by the currently selected checkpoint.\\
Decoded artifacts & Artifacts with at least one stored decode attempt. This
does not mean that every attempt succeeded.\\
Process groups & Number of distinct, nonempty group names.\\
Studio workflow / Open view & Four cards describe the four main views. Each
\ui{Open view} link opens its corresponding page.\\
Representation path & A text diagram showing file bytes, parsed input,
canonical input, latent numbers, comparison, decoded tree, derived net, and
behavior. It describes the pipeline; it is not a control.\\
\end{controls}

In the representation path, $\mu$ (mu) is the embedding's mean, used for
ordinary deterministic comparisons. $\log\sigma^2$ is its log variance,
which describes the spread used in experimental latent sampling.

\section{Artifact workspace}
\label{sec:workspace}
The heading is \ui{Import once. Explore immediately.} Start here to provide
the inputs used by the other views.

\subsection{Import artifacts}
\begin{controls}
Upload your artifacts & Choose one or more XES, PTML, or PNML files. Parsing,
preparation, and encoding start as soon as the files are selected.\\
Bundled artifacts & Choose one repository example or \ui{All 9 artifacts}.
The examples are \code{receipt}, \code{roadtraffic100traces}, and
\code{running-example}, each in all three formats. The displayed total follows
the files available on the server.\\
Import bundled selection & Imports and encodes the selected example or set.
It is disabled while \ui{Choose an example} is selected.\\
Import progress and result message & Shows how many files were processed and
successfully embedded. A warning lists files that failed or reached a limit.\\
\end{controls}

The workspace accepts at most 25 artifacts, with a 50 MiB limit per file.
A repeated filename with identical content reuses its workspace entry and
goes through encoding again. Uploaded files start without a process group;
bundled examples receive their base filename as a suggested group.

\subsection{Advanced preprocessing}
\label{sec:preprocessing}
Preprocessing decides what the encoder receives. Set these options
\emph{before} selecting files or pressing the bundled-import button.
Editing them alone does not re-encode a stored item. Re-import to apply new
encoder limits or sampling settings. If you change parsing settings such as
the activity attribute, remove the existing item, clear inference and evaluation
caches in the Checkpoint dashboard, and then import it again. Duplicate imports
reuse the already parsed artifact, and cached encoder inputs may otherwise
still reflect the earlier parsing choices.

\begin{controls}
Activity attribute & Event field containing the activity name.
Default: \code{concept:name}.\\
Case identifier attribute & Field identifying the case to which an event
belongs. Default: \code{case:concept:name}.\\
Append lifecycle transition to activity labels & Off by default. Distinguishes
lifecycle stages such as an activity's start and completion when the event
contains that information. This can increase the number of distinct labels.\\
Remove empty traces & On by default. Removes cases that contain no activities.\\
Guess missing PNML final marking & Off by default. Allows the parser to try
inferring the final token state if the PNML file does not supply it. Read any
resulting warning.\\
Maximum encoder traces & Largest number of traces passed to the log encoder.
Default 128; allowed range 1--100,000.\\
Maximum trace length & Largest number of events retained in each selected
trace. Longer traces are clipped. Default 128; range 1--10,000.\\
Maximum workspace events per log & Rejects an oversized log for encoding
rather than silently reducing it. Default 1,000,000; range 1--10,000,000.\\
Trace selection & Chooses which traces enter the encoder; the five options
are explained below. Default: \ui{First N}.\\
Compress duplicate traces to variants before selection & Off by default.
Keeps one copy of each distinct sequence before selecting traces. This changes
the frequency information available to selection and encoding.\\
Variant coverage target & Appears for \ui{Variant Coverage}. Default 0.80;
range 0.10--1.00. Sets the fraction of cases represented by the chosen frequent
variants, subject to the trace-count limit.\\
Selection seed & Default 13. Makes random selection or shuffling repeatable
for the same input and settings.\\
Maximum Petri-net nodes & Maximum total places and transitions allowed for
encoding. Default 512; range 1--100,000. Larger nets are rejected, not cut down.\\
Maximum process-tree tokens & Maximum encoded tree length, including boundary
tokens. Default 512; range 2--100,000. Larger trees are rejected, not clipped.\\
\end{controls}

A \emph{variant} is a distinct activity sequence. Many cases can follow the
same variant.
\begin{controls}
First N & Take the first traces in file order, up to the limit.\\
Seeded Random & Take a reproducible random subset without selecting the same
trace position twice.\\
Frequency Preserving & Allocate the trace budget approximately in proportion
to variant frequencies, then shuffle the selected cases using the seed.\\
Most Frequent Variants & Fill the budget from the most frequent variants
first. Repeated copies can be selected.\\
Variant Coverage & Add frequent variants until their source frequencies reach
the coverage target or the trace budget is exhausted. It can return fewer
traces than the maximum.\\
\end{controls}

These settings do not enlarge a checkpoint's activity vocabulary. Tree
operators with too many children may be regrouped into smaller operators
without dropping branches; the preprocessing ledger records this change.

\subsection{Workspace register and selection}
\begin{controls}
Workspace register & One row per artifact. \ui{Artifact}, \ui{Group}, and
\ui{Modality} identify it. \ui{Size} means events for a log, nodes for a tree,
or nodes for a net. \ui{Labels} counts distinct activities for logs/trees,
but visible transitions for nets.\\
State / Encoded / Decoded / Validation & Processing status; whether a stored
embedding is usable; whether a decode has been attempted; and whether the latest
stored decode passed all success checks. A dash means no decode is recorded.\\
Inspect artifact & Chooses the item shown in the panels below. Selecting a
table row is not the inspection control.\\
Process group & Type a shared name for artifacts that describe the same process.
For example, give a log and its tree the group \code{Order process}. A group
is your assertion of a relationship; the app does not prove equivalence.\\
\end{controls}

Typical states are \code{parsed}, \code{canonicalizing}, \code{encoding},
\code{embedding ready}, \code{input limitation}, and \code{encoding failed}.
Decode attempts can add \code{Petri conversion valid}, \code{invalid tree},
or \code{decode limit reached}. Read the detailed warnings and validation
gates to understand a failure rather than relying on the short state alone.

\subsection{Source preview}
This tab shows the parsed source before encoder trace selection and clipping.
For logs, parsing choices such as removing empty traces already apply.

An \ui{Event log} preview contains:
\begin{controls}
Traces / Events / Activities / Variants / Mean length & Case count, event
count, distinct activity count, distinct sequence count, and average events
per trace.\\
Activity frequency / Trace length & Counts of events by activity and counts
of traces by their length.\\
Start activities / End activities & Counts of the first and last activities
of traces.\\
Directly-follows graph & An arrow from A to B means that B immediately follows
A in some trace. Edge labels show counts or percentages. Green activities
can start a trace; orange activities can end one; other activities are purple.
If an activity both starts and ends traces, the start color takes priority.\\
Minimum edge frequency & Hide arrows with smaller counts. Default 1.\\
Maximum displayed edges & Show only the most frequent remaining arrows.
Default 30; range 5--100.\\
Activity search & Keep arrows with a matching source or target activity name;
matching ignores letter case. This filters the drawing, not the uploaded log.\\
Frequency & \ui{Absolute} shows counts. \ui{Relative} shows the percentage
of all directly-follows occurrences in the source, including hidden edges.\\
Most frequent variants & Table of the 100 most frequent sequences at most,
with their counts.\\
\end{controls}

A \ui{Process tree} preview offers \ui{Tree labels}: select
\ui{Original labels} for source names or \ui{Canonical labels} for model
names such as \code{A0}. The tree viewer controls are explained in
Section~\ref{sec:viewers}. A statistics row gives size, depth, activity and
silent leaves, counts of each operator, maximum arity, prefix-token length,
prefix tokens, and distinct activities.

A \ui{Petri net} preview uses the net viewer in Section~\ref{sec:viewers}.
Its statistics include places, visible and invisible transitions, arcs, total
nodes, initial and final token counts, duplicate visible labels, and isolated
nodes (nodes with no incident arcs).

\subsection{Model-facing input}
\begin{controls}
Canonical activity map & Shows original names and the shorter model names.
For example, \code{Receive order} may become \code{A0}. Names are assigned in
first-seen order within each artifact, so \code{A0} need not mean the same
activity in different files.\\
Original labels / Canonical labels / Checkpoint limit / Reversible & Number of
distinct source names, number mapped, vocabulary capacity, and whether this
input has a complete reversible map. PNML reports \ui{no} for reversibility.\\
Mapping table & \ui{Original label}, \ui{Canonical label}, \ui{Frequency},
and \ui{Used by encoder}. A dash or false entry can indicate an unsupported
label or a label omitted by the structural PNML path.\\
Preprocessing ledger & JSON record of selection, trace counts, clipping,
discarded events, and relevant size limits or tree regrouping. It explains
how the input changed.\\
Model input summary & For logs: input shape, selected traces, and canonical
clipped traces. For trees: prefix tokens, token count, and arity. For nets:
node types, edge count, and initial/final markings.\\
Selected-trace table & For a log, \ui{Trace} is the selected-list index;
\ui{Events} and \ui{Selected sequence} describe the trace before clipping;
\ui{Clipped} flags traces beyond the selected length limit. The canonical
traces in the JSON show what was retained.\\
\end{controls}

\textbf{PNML label limitation.} The external Petri-net encoder uses graph
structure, node types, and markings, but does not use visible activity names.
A net can display its original names in the preview while its embedding omits
them. Decoding from PNML therefore uses canonical names and cannot restore the
source's activity names. The interface repeats this warning at relevant steps.

\subsection{Warnings, embedding status, and attention}
\ui{Warnings \& state} shows the current state, followed by parsing,
preprocessing, or encoding messages. A green message means that no such
limitations were detected; it does not certify output behavior.
\ui{Embedding ready} reports the vector dimension and encoding time in seconds.

When a log encoding includes trace attention, \ui{Event-log trace attention}
shows which selected traces receive more weight when the model combines them.
The table lists \ui{Trace}, \ui{Sequence}, \ui{Length},
\ui{Variant frequency}, \ui{Clipped}, and \ui{Attention weight}, sorted by
weight. The scatter plot puts length on the horizontal axis, weight on the
vertical axis, and frequency in the marker size.

\ui{Aggregate attention by trace variant} combines identical sequences,
sums their attention, and shows \ui{Variant}, \ui{Frequency}, \ui{Length},
and \ui{Attention weight}. Frequencies here concern the selected traces.
Attention is a weight on whole traces, not an explanation of individual events.

\subsection{Export and removal}
\begin{controls}
Prepare workspace export & Builds a snapshot of the current workspace. Press
it again after changes to refresh the download.\\
Download workspace ZIP & Appears after preparation. Contains a
\code{workspace.json} manifest, embedding JSON files, stored decode reports,
available PTML and derived PNML outputs, and simulated XES logs for valid
stored trees. The manifest includes configuration, checkpoint metadata,
available evaluations, pairwise cosine scores, and PCA data.\\
Remove selected & Removes the inspected workspace entry immediately. Its
original disk file is not deleted. Re-import the file to add it again.\\
\end{controls}

The ZIP does not contain the original uploaded files, checkpoint weights, or
a restorable Streamlit session. Keep the original inputs separately. The
cached reference gallery and the explorer's interpolation result are not
workspace artifacts and are not included as stored workspace decodes.

\section{Reading the reusable tree and Petri-net viewers}
\label{sec:viewers}
These controls recur in previews, translations, neighbor inspection, and
reference examples. They change the display, not the underlying model.

\subsection{Process-tree viewer}
Read the tree from the root downward. Activity leaves are green boxes;
silent leaves are grey boxes marked $\tau$. Purple operator nodes show their
type and number of children:
\begin{controls}
SEQ & Perform the child processes in order.\\
XOR & Choose one child process.\\
AND & Perform all child processes, allowing parallel or interleaved work.\\
LOOP & Repeat work according to the loop's body and repeat/exit branches.\\
Focus subtree & Show the \ui{Complete tree} or a particular branch. Node paths
such as \code{0/1} identify successive child positions, counted from zero.\\
Displayed depth & Limit how many levels are drawn from the selected root.
Collapsed nodes show how many descendants are hidden. A one-node view shows
\ui{Single-node tree} instead of a slider.\\
Highlight activity & Highlights matching displayed labels in pink. Matching
ignores letter case; it does not remove nodes or change labels.\\
Prefix tokens & Expands the full tree's linear description, even when the
drawing focuses on one branch. \code{<bos>} means beginning, \code{<eos>}
means end, and \code{ARITY_2}, for example, announces two children.\\
\end{controls}

Node tooltips provide a node identifier, subtree size, and depth. Prefix
notation lists an operator before its children; it is another representation
of the same tree.

\subsection{Petri-net viewer}
Circles are places. Orange boxes are visible transitions (activities). Grey
boxes marked $\tau$ are invisible transitions, which route execution without
recording a visible activity. Arrows are arcs. An initial-token symbol and
count mark initially occupied places; a target-like symbol marks a final place
when no initial-token label takes priority. Tooltips provide node names and
marking details.

\begin{controls}
Graph view & \ui{Full graph} draws the net; \ui{Local neighborhood} draws
nodes near a chosen node; \ui{Structural summary only} shows counts without
the drawing.\\
Hide invisible transitions & Hides silent transitions and their incident arcs
from the drawing. It does not merge those arcs into new connections.\\
Largest component only & Keeps the largest connected part, treating arrows
as connections in either direction for this display selection.\\
Neighborhood center / Radius & Appear in \ui{Local neighborhood}. Select a
node and a radius of 1--4 connections. Neighbors are found in either direction;
this is not an execution-reachability test.\\
Places / Visible / Invisible / Arcs & Counts in the selected component or
neighborhood. Hiding invisible transitions changes the drawing but does not
remove them from these counts.\\
\end{controls}

\section{Translation studio}
\label{sec:translation}
The heading is \ui{From latent point to validated model.} At least one
artifact must have a successful embedding for the active checkpoint.

\subsection{Choose a latent source and decode}
\begin{controls}
Latent sources & Select one or more compatible workspace artifacts. A single
source uses its vector; multiple sources combine their vectors. There is no
automatic requirement that selected items share a process group.\\
Maximum decode length & Token budget for generation. Default 256; range
16--1,024, in steps of 16. Tokens include operators and boundary symbols,
so this is not an activity count.\\
Latent mode: Deterministic mean & Uses $\mu$. For several sources the default
is their equal-weight arithmetic mean. The decoder chooses the highest-scoring
allowed token at each step.\\
Use custom fusion weights: experimental & Appears when several sources are
selected. A number field named after each artifact sets its nonnegative weight,
initially 1.0. Weights are normalized automatically; their sum must exceed zero.\\
Latent mode: Stochastic latent samples: experimental & Draws alternative
vectors using the encoded means and spreads, then decodes them. Token selection
still uses the same greedy, grammar-constrained decoder.\\
Sampling seed & Default 13. Controls the repeatable random draws for a fixed
input and configuration.\\
Alternatives & Number of sampled candidates, 1--12; default 4.\\
Collapse duplicate trees & On by default. Keeps only distinct generated token
sequences in the displayed gallery, so fewer alternatives may appear.\\
Decode process tree & Runs generation. Disabled if no source is selected.
The prefix display and progress bar show generated tokens, current step,
grammar state, and number of allowed next choices.\\
\end{controls}

The \ui{Latent source} explanation card describes the default single-source
or equal-weight choice. Custom weights and sampling override those defaults.
Fusion is most interpretable when the inputs describe the same process.
Sampling illustrates possible output variation; the spread is not a calibrated
confidence score for correctness.

If several distinct sampled results remain, the \ui{Stochastic decode gallery}
shows each alternative's tokens, token count, source/seed label, and validity.
\ui{Inspect alternative} chooses the candidate used by the panels and
downloads below. After changing sources or decoding settings, press
\ui{Decode process tree} again so those panels refer to the new selection.

\subsection{Validation gates}
Each badge checks a different part of the result. Green means the check passed;
red means it failed. Warnings and errors appear below the badges.
\begin{controls}
EOS emitted & The decoder produced the explicit end token \code{<eos>}.\\
Grammar valid & The generated tokens can be read as a complete process tree.\\
Arity valid & Tree operators respect the allowed number of children.\\
Vocabulary valid & Generated token identifiers belong to the model's vocabulary.\\
Petri convertible & The decoded tree was successfully converted to a Petri net.\\
Length limit & Green means the generation limit was \emph{not} reached without
an end token. Red means the token budget was exhausted without EOS.\\
\end{controls}

These are structural and termination checks. They do not show that the tree
reproduces the input behavior or that activity names were fully restored.

\subsection{Decoder progress}
The colored token ribbon distinguishes operators (purple), arity tokens
(blue), activities (green), silent \code{TAU} tokens (grey), and boundary
tokens (orange). The table explains each generation step:
\begin{controls}
Step / Grammar state & Step number and the kind of token the partial tree
currently allows.\\
Chosen / Score & Selected token and its probability among allowed choices at
that step. This is not the probability that the entire discovered process is
correct.\\
Valid choices / Open slots & Number of allowed next tokens and the child
positions that still need to be filled.\\
Subtree position / EOS & Where the next tree content belongs, and whether
the end token was emitted at that step.\\
Top valid & Highest-scoring allowed token choices, with probabilities
normalized over the allowed set.\\
Masked high scores & Highly scored candidates blocked by the grammar. They
were not chosen because they would violate the current syntax. Their displayed
scores come from the full vocabulary before masking.\\
\end{controls}

\subsection{Decoded process tree and Derived Petri net}
\ui{Decoded process tree} has two subtabs. \ui{Canonical labels} shows model
names such as \code{A0}. \ui{Original labels} uses the source map when one is
available. For fusion, restoration requires identical original-to-canonical
maps in all sources. A generated activity absent from the map stays canonical.
If any source is PNML, original-name restoration is unavailable. Thus an
\ui{Original labels} tab can still show canonical names; read the warnings.

\ui{Derived Petri net} shows the deterministic conversion of the canonical
decoded tree. It is not a separately generated neural output. The reusable net
viewer is followed by full-net counts of \ui{Places}, \ui{Transitions},
\ui{Silent}, and \ui{Arcs}. These extra counts describe the full result even
if the viewer above is filtered.

\subsection{Source comparison}
This tab compares the displayed decode with each source used in its generation.
A summary table contains scalar measurements; additional panels depend on the
source type.

\begin{controls}
Process-tree source & Side-by-side source and decoded trees. Exact tree/prefix
matches test structural representations. Token edit distance counts insertions,
deletions, and replacements; its normalized version divides by the longer
token sequence. Lower is closer. Size, depth, operator/leaf overlaps, and count
differences describe other structural changes. The overlap scores compare
counts, not execution behavior.\\
Event-log source & Compares the log with traces simulated from the decoded
tree. Charts show activity frequencies and trace lengths; paired tables show
the top 20 variants and top 30 directly-follows relations. The summary includes
distribution distances and simulated trace count.\\
Petri-net source & Side-by-side source and derived nets, with source/decoded
counts for places, transitions, visible/invisible transitions, arcs, initial
and final tokens, and duplicate visible labels. Label preservation is explicitly
not evaluated.\\
\end{controls}

Distribution distances describe differences between observed frequency
patterns. An L1 distance of zero means identical compared distributions;
smaller is closer. The fields \code{variant_l1}, \code{directly_follows_l1},
and \code{length_l1} compare variants, directly-follows pairs, and trace
lengths; \code{mean_l1} averages those three. The fields
\code{activity_frequency_l1}, \code{start_activity_l1}, and
\code{end_activity_l1} compare activities and trace starts/ends.
\code{aggregate_distribution_l1} averages the first three distances plus
activity-frequency distance. \code{mean_trace_length_difference} is the
absolute difference in average events per trace. Simulated samples are only
samples of the model's behavior; equal summaries do not prove process
equivalence.

\subsection{Export decoded model}
\begin{controls}
Restore original labels in PTML & On by default where allowed. Affects the
PTML download only. Hidden for PNML-derived output. It cannot restore a name
that has no source mapping.\\
Download PTML & Saves the tree. Offered when grammar validation succeeds.\\
Download validation JSON & Saves source identifiers, tokens, trees, label map,
validation flags, warnings, errors, decode time, and decoder settings.
Also requires a grammar-valid tree in this interface.\\
Download derived PNML & Saves the converted net when conversion succeeded.
It retains canonical tree labels; the PTML restoration checkbox does not
relabel this net.\\
Download simulated XES & Saves a generated event log from the decoded tree,
using 100 traces and seed 13. Uses restored labels when a restored tree is
available. Requires a grammar-valid tree.\\
\end{controls}

A download button's presence is not a substitute for checking every validation
gate. The report records the distinction between a parsed tree and a fully
successful decode.

\section{Latent explorer}
\label{sec:latent}
The heading is \ui{Make multimodal agreement visible.} The page uses
nonempty, error-free embeddings from the active checkpoint. Different trained
models can assign different meanings to the same vector coordinates, even
when their dimensions match.

\subsection{Embedding register}
\begin{controls}
Artifact / Group / Modality & File name, assigned process group, and input type.\\
Dimension / Norm & Number of coordinates and vector length. A larger norm
does not mean a better process or a better prediction.\\
Latent spread & Average latent standard deviation. Describes sampling spread,
not measured prediction accuracy.\\
Nearest artifact / Similarity & Another workspace artifact with the highest
cosine similarity, and that score. Empty if there is no other artifact.\\
Encoded at / Checkpoint & Encoding timestamp and saved-model identifier.\\
\end{controls}

\subsection{PCA projection}
PCA reduces the vectors to two plotted coordinates. Each point is an embedding;
color and shape identify its modality. Grey lines join points with the same
nonempty group. Hover for the artifact, group, modality, and coordinates;
the plot supports interactive panning and zooming.

\begin{controls}
Show equal-weight fused group means & On by default. Adds a mean point for
each group having at least two modalities, using the first artifact of each
modality in that group.\\
Include cached reference gallery & Adds previously generated synthetic
examples. These are labelled as references and use the same checkpoint.\\
Principal component 1 / Principal component 2 & Horizontal and vertical
projection coordinates, not named process properties.\\
Explained variance & Percentage of the vector variation retained by each
axis. The two percentages need not add to 100\%.\\
\end{controls}

At least three points with variation in two independent directions are needed
for the plot. Otherwise the panel recommends pairwise distances. Adding
references or fused means recomputes the projection, so existing points can
move on screen without their embeddings changing. Nearby projected points
are a useful clue, not proof that the original processes are equivalent.

\subsection{Similarity heatmap}
\begin{controls}
Ordering & Choose \ui{Upload order}, \ui{Process group}, \ui{Modality},
\ui{Similarity to selected}, or \ui{Hierarchical clustering}. This rearranges
rows and columns without changing their values. Clustering can fall back to
upload order if unavailable.\\
Similarity anchor & Appears for \ui{Similarity to selected}. Choose the
artifact whose most similar neighbors should appear first.\\
Pairwise measure: Cosine similarity & A colored grid comparing vector
directions. Values near 1 mean similar directions, 0 roughly perpendicular,
and $-1$ opposite. Green is high, dark is near zero, red is low. Hover a cell
for both names and its score. This is not an accuracy percentage.\\
Pairwise measure: Euclidean distance & A numeric table of straight-line
distances in the full latent space. Zero means identical vectors; smaller is
closer.\\
Pairwise measure: Normalized Euclidean distance & Euclidean distance divided
by the square root of the latent dimension. It is not constrained to 0--1.\\
\end{controls}

This panel compares the workspace embeddings; projection reference points
and fused group points are not added to this matrix.

\subsection{Group agreement}
This table checks whether different representations assigned to one process
group lie close together. It uses the first artifact of each modality in the
group and requires at least two different modalities. Duplicating a log within
a group therefore does not create a second modality for this calculation.

\begin{controls}
Process group / Artifacts & Group name and the number of modality
representatives actually compared.\\
Mean / Maximum within-group distance & Average and largest pairwise
Euclidean distances. Smaller means closer agreement in the learned space.\\
Tree--Log / Tree--Petri / Log--Petri cosine & Pairwise cosine scores for the
modalities present in the group. Larger means closer directional agreement.\\
Distance from fused & Each modality's Euclidean distance from the equal-weight
group mean. A larger value identifies a representation farther from that mean.\\
Rank groups by & \ui{Best agreement} sorts by smallest mean distance;
\ui{Worst agreement} reverses it. \ui{Largest Petri discrepancy} and
\ui{Largest event-log discrepancy} sort by the respective distance from the
fused mean.\\
\end{controls}

\subsection{Nearest neighbors}
\begin{controls}
Neighbor pool & \ui{Workspace} or \ui{Workspace + reference gallery}.
The latter requires a built reference gallery.\\
Selected artifact & Workspace artifact for which to find neighbors.\\
Modality filter & \ui{Any modality}, \ui{Other modalities only}, or
\ui{Same modality only}.\\
Neighbor table & Other artifacts sorted by decreasing cosine similarity,
then increasing Euclidean distance. Includes identifiers, names, modality,
both measures, \code{same_process_group}, and \code{source}
(workspace or reference).\\
behavioral\_distance & Appears when traces are available for both entries,
either directly or through an event log in the same process group. Uses
\code{mean_l1}; smaller means closer trace-distribution summaries. It does
not determine the neighbor ordering.\\
Inspect neighbor source & Select one of the first 20 neighbors to see its
tree, net, or top 20 log variants. A workspace neighbor's latest decoded tree
also appears when available. A reference neighbor shows its generating tree.\\
\end{controls}

\subsection{Dimensions}
\ui{Compare artifacts} selects vectors for coordinate-by-coordinate line
charts; the first three artifacts are selected initially. \ui{Latent mean by
dimension} plots $\mu$. \ui{Latent standard deviation by dimension} plots
$\exp(\tfrac12\text{logvar})$. With several selected artifacts,
\ui{Difference from selected fused mean} plots each vector minus their
arithmetic mean. Here every selected artifact contributes. Coordinates are
intentionally unnamed: a peak is not labelled as a particular process property.

\subsection{Interpolation}
This is an exploratory view requiring two encoded artifacts.
\begin{controls}
Left endpoint / Right endpoint & Choose two different workspace vectors.\\
alpha & Slider from 0 to 1 in steps of 0.05; default 0.5. Zero uses the left
vector, one uses the right, and 0.5 takes their midpoint.\\
Decode interpolation point & Decodes $(1-\alpha)\mu_{\mathrm{left}}+
\alpha\mu_{\mathrm{right}}$ with a fixed 256-token budget.\\
Result & Generated token sequence, \ui{Tree size}, \ui{Decode tokens}, and
\ui{Petri conversion} status, followed by the available tree and net viewers.\\
\end{controls}

Intermediate vectors need not produce a smooth or meaningful sequence of
process changes. This view does not restore original labels. Press the decode
button again after changing an endpoint or alpha; the stored result otherwise
remains the previous one.

\subsection{Reference gallery}
This tab creates synthetic comparison examples for a small workspace. It is
available after at least one workspace artifact is encoded.
\begin{controls}
Reference processes & Number of synthetic processes, 3--100; default 12.\\
Reference seed & Default 13. Controls example generation.\\
Traces per process & Number of generated log traces per process, 8--512;
default 64.\\
Build reference gallery & Generates processes and encodes each tree, log,
and net with the active checkpoint. A progress bar counts completed processes.\\
Cached-gallery message and table & Shows process and embedding totals, then
\ui{Reference}, \ui{Group}, \ui{Modality}, \ui{Tree size}, \ui{Traces},
and \ui{Petri nodes}. Each process contributes three modality embeddings.\\
Preview reference process & Select a reference group. \ui{Process tree}
shows its tree; \ui{Event-log behavior} shows up to 30 frequent variants;
\ui{Petri structure} shows its net.\\
\end{controls}

References live in a checkpoint-specific session cache. They can be included
in PCA and nearest-neighbor searches. Building the gallery again replaces it;
switching checkpoints invalidates the old gallery when this page is visited.

\subsection{Inspect and export one vector}
\begin{controls}
Embedding & Select a workspace vector to inspect and download.\\
Raw latent mean and log variance & Expands a table of \ui{Dimension},
\code{mu}, and \code{logvar}. Dimension indices start at zero.\\
JSON & Full encoding record, including vectors, source/preprocessing metadata,
checkpoint identifier, canonical map, runtime, and diagnostics.\\
CSV & One row per dimension with mean, log variance, and repeated metadata.
Useful for spreadsheet analysis.\\
NumPy & A \code{.npz} archive containing \code{mu}, \code{logvar}, and
JSON metadata for numerical analysis.\\
\end{controls}

\section{Checkpoint dashboard}
\label{sec:checkpoint}
The heading is \ui{Inspect the run behind the latent space.} This page reads
saved training information; it has no training-start button.

\subsection{Model facts}
\begin{controls}
Epoch / Best epoch & The checkpoint's saved epoch and the latest epoch marked
as best in its embedded history. An epoch is one training pass through the
data. A missing value is shown as a dash.\\
Latent dimension / Hidden dimension & Size of the shared vectors and the
model's internal representations.\\
Activity limit / Best validation loss & Maximum supported activity vocabulary
and best recorded validation objective. Lower loss is better for the same
objective and comparable data.\\
early stopped & Annotation inferred from the recorded patience and final
non-improving-epoch count, when that information is present.\\
Facts table & \ui{Tree arity}, \ui{Tree token limit}, \ui{Petri node limit},
\ui{Training traces/sample}, \ui{Trace length limit}, \ui{Checkpoint type},
and \ui{Training/file timestamp}. These are recorded model/run facts; compare
them with the current import settings rather than treating every value as a
live workspace setting.\\
\end{controls}

\subsection{Training history}
The main line chart compares training and validation loss against epoch.
Training loss concerns the examples used to learn; validation loss concerns
held-out examples. \ui{Validation component losses} separates the terms
defined below. \ui{Generalization gap} plots validation loss minus training
loss. A growing positive gap can indicate that performance on training data
is improving faster than performance on validation data.

\begin{controls}
tree\_reconstruction & Error when predicting a tree from its own representation.\\
trace\_to\_tree / petri\_to\_tree & Errors when predicting tree tokens from
a log or net representation.\\
latent\_alignment & Penalty for separating matching representations.\\
contrastive & Encourages matching behavior representations to be closer than
unmatched ones.\\
kl & Regularization of the latent distributions, not a behavior-accuracy score.\\
loss & Weighted total of the component objectives. Component scales can differ.\\
\end{controls}

The history table includes epoch, learning rate, epoch duration in seconds,
best-epoch flag, best validation loss, training/validation components, and
available gaps. If embedded history is absent, the page tries the adjacent
server-side \code{training_metrics.csv}; it shows its table and available
total-loss curves. If neither exists, an information message explains the
absence. A usable checkpoint does not need a complete history.

\subsection{Configuration}
\ui{Training configuration} and \ui{Synthetic-data configuration} show saved
settings as JSON. \ui{Complete checkpoint metadata} includes all recorded
facts. These displays are read-only.

\ui{Stage caches} lists \ui{Stage}, \ui{Entries}, and \ui{Invalidated by}.
A cache is temporary storage that avoids repeating calculations. The stages
are \ui{Parsed Artifact}, \ui{Preprocessed Input}, \ui{Embedding},
\ui{Decode}, \ui{Simulation}, and \ui{Evaluation}. The third column describes
which inputs identify reusable work, such as a file's content and parse
settings, or a latent vector and decoder configuration.

\ui{Clear inference and evaluation caches} clears the preprocessed-input,
embedding, decode, simulation, and evaluation caches. It preserves the parsed
artifact cache and the workspace items. It does not erase already displayed
results or force an automatic re-encoding; repeat the relevant action when
you want a new result.

\subsection{Compare checkpoints}
This table reads all trusted checkpoints and compares \ui{Checkpoint},
\ui{Type}, \ui{Epoch}, \ui{Latent}, \ui{Hidden}, \ui{Activities},
\ui{Arity}, and \ui{Best validation loss}. A file that cannot load receives an
\ui{Error} entry. This is a metadata comparison; it does not re-evaluate all
models on your workspace. Select another active model through the sidebar.

\section{Common situations and next steps}
\begin{controls}
No checkpoint found / load failed & A loadable trusted \code{.pt} file must
be installed in the configured server directory. Browser artifact upload does
not accept checkpoints.\\
No artifacts or no embeddings & Import an example in Artifact workspace.
If parsing succeeded but encoding did not, inspect \ui{Warnings \& state}.\\
Too many activities & The checkpoint vocabulary is too small for the input.
Use a compatible checkpoint or an appropriate smaller input; increasing trace
or node limits does not expand the vocabulary.\\
Trace selection or clipping warning & Inspect the ledger and canonical traces.
Increase the relevant limits or change selection if needed, then re-import.\\
Tree/net size limit & Increase the corresponding preprocessing limit when
appropriate, then re-import. Those structures are not silently truncated.\\
Checkpoint mismatch & Re-import affected items with the active checkpoint.
Use the same checkpoint for all vectors being compared or exported.\\
No PCA plot & Add more varied artifacts or build references. Use the pairwise
matrix when fewer than three independent points are available.\\
No group agreement & Assign a shared nonempty group to at least two different
encoded modalities that describe the same process.\\
No EOS / incomplete decode & Inspect the token trace and error messages. A
larger decode budget may help; it does not guarantee a useful model.\\
Original labels are missing & Check the canonical map. PNML-derived output
has no restoration; fused sources need identical maps; newly generated labels
may have no source counterpart.\\
A result seems unchanged & Settings do not replace stored action results.
Press decode, build, run, or prepare again, as appropriate. Re-prepare a ZIP
after workspace changes.\\
Evaluation page is missing & It is not registered in the current navigation.
Appendix~\ref{sec:evaluation} describes the separate source page.\\
\end{controls}

\appendix
\section{Quality evaluation: page present in the source only}
\label{sec:evaluation}
\code{pages/04_evaluation.py} defines a page titled \ui{Quality evaluation},
with the heading \ui{Results arrive while the run is still alive.} The normal
\code{streamlit_app.py} navigation does not register it. Its ``View 04''
heading should not be confused with the visible Checkpoint dashboard. The
controls below describe that source page if it is integrated into a deployment;
this guide does not change the navigation.

The page requires workspace embeddings for the active checkpoint. It offers
five tabs whose result rows arrive progressively, so completed work can be
inspected before the entire evaluation finishes.

\subsection{Shared evaluation controls}
\begin{controls}
Evaluation artifacts & Multiselect of workspace artifacts; initially all are
selected. Applies to Decode \& behavior, Discovery comparison, and Embedding
baselines. Cross-modal retrieval uses all eligible grouped workspace items;
Neural losses reads a server-side data split.\\
Decode token limit & Default 256; range 16--1,024 in steps of 16. Used by
decode evaluation and discovery comparison.\\
Quick simulation traces & Default 100; range 10--10,000, in steps of 10.
Number of simulated traces for quick decode/behavior comparisons.\\
\end{controls}

\subsection{Decode \& behavior}
\begin{controls}
Also run expensive alignment fitness and precision & Off by default.
Adds alignment-based comparison where an event-log source is available.\\
Run progressive evaluation & Evaluates the selected artifacts and equal-weight
fused groups with at least two modalities, using the first item of each
modality. The progress bar shows completed tasks and elapsed time.\\
Result table & Contains artifact, modality, latent source, EOS, tree validity,
Petri conversion, decode length/time, source-comparison values, and any
available alignment scores or errors. Comparison fields depend on input type
as described in Section~\ref{sec:translation}.\\
Rate chart & Fractions of evaluated rows with EOS, a valid tree, and successful
Petri conversion. A rate of 1 means every evaluated row passed that check.\\
Failed example & Appears for workspace examples whose tree or conversion
failed. Select one for closer inspection.\\
Open failed example in Translation studio & Navigates to translation with
that artifact preselected so you can decode and inspect it there. Requires
that page to be registered in the deployment.\\
Download evaluation JSON / Download evaluation CSV & Saves this tab's
completed rows as structured text or a spreadsheet-friendly table.\\
\end{controls}

\emph{Fitness} measures how well the model can reproduce observed behavior.
\emph{Precision} measures how much additional behavior the model allows beyond
the observations. F1 combines them. Higher is better for these measures, but
they answer different questions from structural validation or quick
distribution distances. Alignment runs can take considerably longer.

\subsection{Neural losses}
This tab evaluates the model's learning objectives on stored synthetic data.
It does not calculate these losses on the selected uploaded files.
\begin{controls}
Server-side split & Choose \ui{test}, \ui{validation}, or \ui{training};
default \ui{test}. The caption shows the resolved sample file.\\
Batch size & Samples processed together, 1--512; default 16.\\
Maximum batches & Maximum number of batches to evaluate, 1--10,000;
default 20. This may evaluate only part of the split.\\
Use deterministic latent means & On by default. Uses means instead of sampled
latent vectors.\\
Neural evaluation seed & Default 13. Controls repeatable random evaluation.\\
Run neural-loss evaluation & Processes batches, updates a running-loss chart,
and shows batch and running-average measurements in a table. Loss components
are explained in Section~\ref{sec:checkpoint}.\\
\end{controls}

The sample path is \code{data/<split>/samples.jsonl}, or uses the root named
by \code{PROC_ROSETTA_DATA_DIR}. If it is missing, the page displays an
information message instead of the run button.

\subsection{Cross-modal retrieval}
\ui{Run cross-modal retrieval} checks all six directions between tree, log,
and net. For example, tree-to-log asks whether a tree's nearest log belongs to
the same process group. Each direction uses groups with both required
modalities. If a group has several artifacts of one modality, this evaluator
uses the last one for that modality; use one per modality for a clear comparison.

\begin{controls}
direction / count / available & Modality direction, number of paired groups,
and whether a comparison can be made. A reason explains missing results.\\
top1\_accuracy / top3\_accuracy / top5\_accuracy & Fraction of queries whose
matching group is among the first 1, 3, or 5 candidates. Larger is better.\\
mrr & Mean reciprocal rank: a first-place match contributes 1, second place
contributes $1/2$, and so on. Larger is better.\\
mean\_rank / median\_rank & Average and middle matching positions. Smaller
is better; rank 1 is best.\\
rank\_histogram & Counts of matches at each rank, included in the result table.\\
Retrieval bar chart & Compares the top-1, top-3, top-5, and MRR values across
available directions.\\
\end{controls}

With very few paired groups, a high top-$k$ score can be trivial. Use the
\code{count} field when interpreting it.

\subsection{Discovery comparison}
\begin{controls}
Run alignment fitness and precision & Off by default. Enables the more
expensive conformance measurements for both discovery methods.\\
Compare with Inductive Miner & Uses selected event logs to generate one model
with ProcRosetta and one with the classical Inductive Miner method.\\
Comparison table & Artifact, method, model-discovery and conversion status,
runtime, places, transitions, silent transitions, arcs, errors, and available
fitness, precision, and F1. Without alignment enabled, those scores remain
empty; they are not zero.\\
F1 chart & Appears when usable F1 values exist and compares the two methods
for each artifact.\\
\end{controls}

The methods optimize different objectives. Runtime fields also cover different
steps: the ProcRosetta row measures its decode, while the Inductive Miner row
includes its discovery work and any requested conformance calculation. Do not
treat them as a controlled end-to-end speed benchmark.

\subsection{Embedding baselines}
This tab compares learned vectors with simpler descriptions of logs or nets.
\begin{controls}
Include optional PM4Py Petri Node2Vec & Off by default. Adds a graph-embedding
method based on walks through the net. Requires at least two process groups
with paired XES and PNML inputs.\\
Node2Vec dimensions / Walks per node / Node2Vec epochs & Appear when that
option is enabled. Defaults 64, 5, and 5; ranges 4--512, 1--100, and 1--100.
They control vector size and the amount of graph-embedding training.\\
Run embedding baselines & Runs available methods progressively, with a progress
bar, full JSON results, a ranked summary table, and available score charts.\\
\end{controls}

The methods include activity counts, variant distributions, directly-follows
pairs, eventually-follows pairs (one activity occurs later than another), PM4Py
log-case features, and ProcRosetta event-log means. Additional methods use
Petri structural counts and ProcRosetta tree, net, and fused means aligned
through process groups. Log comparisons need at least two logs; the grouped
methods need at least two groups containing logs and the required representation.
Unavailable methods report a reason.

\begin{controls}
Method / Dimension / Runtime & Name of the representation, vector size, and
available runtime in seconds. A missing runtime is not a zero-time run.\\
Behavior Spearman / Behavior Pearson & Correlations between embedding
distances and trace-behavior distances. Spearman compares their order;
Pearson compares their linear association. Larger positive values indicate
stronger agreement for these samples.\\
Nearest-neighbor behavior & Mean behavior distance to the nearest embedded
neighbor. Smaller means more behaviorally similar neighbors.\\
Improvement over random & Improvement relative to random-neighbor behavior
distance. Larger positive values indicate an advantage over that reference.\\
Top-1 neighbor agreement / Top-3 neighbor agreement & Fraction of items whose
closest behavioral neighbor appears first, or among the first three, in the
embedding's neighbor ordering.\\
Rank & Orders methods by Behavior Spearman, largest first, with unavailable
values last. The chart shows available Spearman and nearest-neighbor behavior
values, whose scales and preferred directions differ.\\
\end{controls}

Small sets or constant distances can make correlations uninformative or
unavailable. These comparisons measure the selected samples and their
frequency-based behavior summaries, not universal process quality.

\subsection{Complete evaluation report and cached results}
\ui{Download complete JSON report} saves the checkpoint identifier, current
shared configuration, and the stored results of all five tabs. Tabs that have
not run contribute empty lists. Unlike the per-tab CSV, this includes the full
collection of evaluation sections.

An evaluation may reuse cached results, with an information message explaining
the reuse. Changing settings does not remove old displayed rows, and the
complete report collects the stored results. Rerun the relevant tabs after
changing configuration or checkpoint before treating the report as one
consistent experiment. The dashboard's cache-clear button does not erase
those already stored result lists.

\section{Document sources and compilation}
\label{sec:sources}
The guide was checked against repository revision \code{a1ea7ff}. The following
files determine the interface described here; the application source takes
precedence over older screenshots or descriptions.

\begin{controls}
Navigation and overview & \path{streamlit_app.py};
\path{proc_rosetta_ui/home_page.py}.\\
Shared settings and state & \path{proc_rosetta_ui/common.py};
\path{proc_rosetta_ui/app_state.py}; \path{proc_rosetta_ui/cache_service.py}.\\
Main pages & \path{pages/01_workspace.py}; \path{pages/02_translation.py};
\path{pages/03_latent_explorer.py}; \path{pages/05_checkpoint.py}.\\
Separate evaluation page & \path{pages/04_evaluation.py};
\path{src/proc_rosetta/evaluation_iterators.py}.\\
Shared viewers & The Python modules under
\path{proc_rosetta_ui/visualizations/}; \path{src/proc_rosetta/visualization_data.py}.\\
Input, decoding, and exports & \path{src/proc_rosetta/artifact_io.py};
\path{src/proc_rosetta/inference.py}; \path{proc_rosetta_ui/ui_types.py};
\path{proc_rosetta_ui/workspace_service.py};
\path{proc_rosetta_ui/encoding_service.py};
\path{proc_rosetta_ui/decoding_service.py};
\path{proc_rosetta_ui/export_service.py}.\\
Training and behavior measurements & \path{src/proc_rosetta/training.py};
\path{src/proc_rosetta/losses.py}; \path{src/proc_rosetta/behavior.py};
\path{src/proc_rosetta/benchmarks.py}.\\
\end{controls}

Build this standalone guide from the repository root:
\begin{verbatim}
latexmk -pdf -interaction=nonstopmode -halt-on-error \
  -outdir=paper paper/streamlit_interface_guide.tex
\end{verbatim}
The output is \path{paper/streamlit_interface_guide.pdf}. A normal TeX Live
installation with the packages listed in this file's preamble is sufficient;
no bibliography, external images, main-paper files, or shell-escape option is
required.

\end{document}
